% =============================================================================
%  Fragment: Upper Volume Articles 67–73; Upper/Lower Volume transition;
%            Lower Volume Articles 1–2
%  Source:   translation/pages_59-62.md
% =============================================================================

% --- Note: Chapter 66 commentary conclusion is carried over from the
%     previous fragment. It is included here to preserve continuity. ---

% --- Chapter 66 Commentary (conclusion, carried from previous page) ---

For example, suppose one believes this rice will generally rise but judges
that the advance will be delayed somewhat --- that such-and-such bearish
factors remain, and that the price will dip in the near term. One thinks
to sell a little, intending to scoop up the dip. But this should never be
done. Invariably, just when one expects the decline, the market rises
instead. People are stubborn creatures: when one has committed to a sell
position out of firm conviction that prices must fall, and then the market
rises against that conviction, one finds it all the harder to buy.
Wavering and procrastinating, one falls behind the advance, and in the
end, seized by frustration, one ends up selling in the very month one
should have been buying --- and suffers enormous losses. This calls for
the greatest caution. If one expects prices to rise next month, one
should simply wait for a buying opportunity on the dip.

% =============================================================================
%  Article 67
% =============================================================================
\article{67}{Trading on Personal Notion Alone Is Useless}

\begin{sokyubox}
{\small\itshape\color{sokyubar} \jp{第六十七章\quad 一分の存念にて賣買無用の事}}

\medskip
One must never trade on personal notion alone. Place the Three-Position
Tradition (\jt{三位の傳}{sanmi no den}) at the fore, fix one's stance
(\jt{立羽}{tateba}), and commit decisively --- whether buying or selling.
\end{sokyubox}

\commentary
The original text says one must not trade on personal notion. Taken at
face value, this sounds as though one must never trade on one's own
judgment, and one feels a strong urge to object. But the key lies in the
three characters ``personal notion'' (\jp{我一分の}, \textit{ware ichibun
no}). What S\=oky\=u warns against is trading carelessly, without having
assessed the position of the ceiling and floor, and without having
settled on a buying stance or a selling stance. His meaning is this:
combine the Heavenly Stems (\jt{十干}{jikkan}) of the months, consult
the charts of monthly and bale-count position-reckoning to grasp the
highs and lows, and only after a fixed stance --- sell if selling, buy if
buying --- has been firmly established should one begin to put on a trade.

% =============================================================================
%  Article 68
% =============================================================================
\article{68}{Following the Market's Swings Without a Fixed Stance}

\begin{sokyubox}
{\small\itshape\color{sokyubar} \jp{第六十八章\quad 立羽もなく高下に連れ候事}}

\medskip
When one follows the market's swings without a fixed stance, trading on
the thought ``it will rise'' or ``it will fall'' --- after five or six
days, as rice moves, one turns bearish and over-sells; after fourteen or
fifteen days, one turns bullish again --- and each time, losses
accumulate. This happens because one rushes the trade, determined to
capture every move. Abandon greed. Ascertain the ceiling
(\jt{天井}{tenj\=o}) and the floor (\jt{底}{soko}). If it is an
up-market, trade up; if a down-market, trade down. Fix one's stance,
check it against the Three-Position Tradition, and commit decisively to
either buying or selling.
\end{sokyubox}

\commentary
This chapter makes plain the reason why trading without a fixed stance
ends in loss. When one has no settled stance --- neither a buying stance
nor a selling stance --- one turns seller whenever the market looks weak,
reverses position (\textit{doten}), and when prices rise a little,
reverses again to the buy side. At each settlement, the result is a loss.
This happens chiefly because one is in a hurry to win. Such trading
amounts to forcing a view where none exists, and loss is inevitable.
Therefore, abandon the desire for gain, do not try to capture every
up-swing and down-swing and range-bound oscillation without exception,
but establish a fixed stance --- buy or sell --- and devote one's efforts
to capturing the move in one direction.

% =============================================================================
%  Article 69
% =============================================================================
\article{69}{Tempering Oneself Through the Three-Position Tradition}

\begin{sokyubox}
{\small\itshape\color{sokyubar} \jp{第六十九章\quad 三位の傳を以て上下鍛練の事}}

\medskip
In this trade, use the Three-Position Tradition to discern the overall
direction of the market. Within a given move, determine how far the
decline extends, how far the advance reaches, and at what level it halts.

At that time, consider the crop conditions in the Kamigata and local
regions, attend diligently to how matters develop from start to finish.
For example, when one has taken a buy position, pay no heed to the
erratic swings in between; hold firm and do not abandon one's stance.
When one's assessment proves correct, one takes profit exactly as
calculated. The approach of ``buy low, sell high'' may sound sensible,
but it is not good practice. When one judges the market to be rising,
consider the whole picture and commit to buying only (\jt{片買}{katagai}).
If one's judgment proves wrong, promptly sell back, withdraw, and observe
the movement of rice carefully. At such times, because rice appears weak,
one is tempted to over-sell --- but this is most inadvisable. Even if one
is utterly convinced that prices must fall, one should not sell. Bear
this in mind.
\end{sokyubox}

\commentary
To observe the market's swings, one must empty the mind and achieve calm
detachment. First, determine whether this is fundamentally a rising
market, a falling market, a range-bound market (\jt{通相塲}{kayoi
s\=oba}), or a consolidation (\jt{保合}{mochai}). Next, in a rising
market, how far does it rise and how far does it pull back? In a falling
market, how far does it fall and how much does it rally? How far does it
oscillate, and at what level does it halt? Investigate the proportions of
the swings. Then factor in crop conditions, stocks of old rice, volume of
supply, the amount of arbitrage rice (\jt{乘米}{nosei-mai}) at each
location, the direction of market sentiment (\jt{人氣}{ninki}), and so
forth. Consider how, in the end, this market will move.

Once one has established a buying stance or a selling stance and has
begun putting on a position in one direction, one must press forward with
that stance to the very end. One must never fix one's eye upon the
intermediate range-bound oscillations. Consider the whole of the market
and commit to buying only or selling only (\jp{片買又片賣},
\textit{katagai mata katauri}). If one's judgment proves wrong, cut
losses promptly, cease trading, rest and calm the mind, observe the
market's course, and only then resume trading.

Furthermore, when one encounters an intermediate rally or dip, one is
tempted to take a small profit --- but until the full measure of the
months and the price level has matured, one should never take profits.

% =============================================================================
%  Article 70
% =============================================================================
\article{70}{Guarding Against Losing Sight of the Market's Overall Direction}

\begin{sokyubox}
{\small\itshape\color{sokyubar} \jp{第七十章\quad 高下に連れ米の一體を失ふ用心の事}}

\medskip
Even one who has thoroughly studied this book will, when carried along by
the market's swings, lose sight of the overall direction of rice. From
time to time, check against the Three-Position Tradition. Weigh whether
rice overall is heading up or down. When one calculates the highs and
lows day after day, month after month, and whenever the spirit stirs, one
attaches a view to every daily swing and trades all year round --- and
when the trade is settled, the result is invariably a loss. In any given
year, there are only two occasions, at most, when one can trade with real
confidence. Use the Three-Position Tradition to consider whether rice
will rise or fall over the next two or three months; then commit
single-mindedly to either buying or selling. If even this feels
uncertain, wait for months on end and consider the moment when the
opportunity aligns. I say again and again: never depart from the
Three-Position Tradition and trade on personal notion alone. This is of
the utmost importance. If one trades on personal notion, loss is beyond
doubt.
\end{sokyubox}

\commentary
However well one has mastered the meaning of each chapter explained in
this book, once one actually opens a position and is buffeted by the
swings, one's habitual discipline vanishes --- one's eye is caught by the
movement, and one can no longer see the larger trend. Therefore, from
time to time, one should practice something like seated meditation
(\textit{zazen}): still the agitation of one's own mind, and based on
the price-tracking charts (\jt{足取表}{ashitori-hy\=o}) --- that is, the
line charts (\jt{罫線表}{keisen-hy\=o}) --- study the proportions of the
swings, the passage of months and days, the direction and intensity of
market sentiment, whether the market is ripe or unripe, and so forth.
Only when one's conviction in the direction is free of doubt should one
open a position. And, as stated in the two preceding chapters, such
moments of assured opportunity occur only two or three times a year. This
is reiterated here with the same meaning, as earnest and repeated
counsel.

% =============================================================================
%  Article 71
% =============================================================================
\article{71}{Selling Rice That Has Risen Before November with a Margin Call Is Useless}

\begin{sokyubox}
{\small\itshape\color{sokyubar} \jp{第七十一章\quad 霜月前引上落引の米賣方無用の事}}

\medskip
Before the November contract (\jt{霜月限}{shimotsuki-giri}), if rice has
risen only fifty or sixty tawara and triggers a margin call
(\jt{落引}{ochibiki}), one should absolutely not sell. If rice has risen
approximately one hundred tawara and triggers a margin call, one may
regard it as a selling opportunity. However, in the case of spring or
summer rice, when arbitrage rice is plentiful, even a hundred-tawara rise
may not contradict the calculation.
\end{sokyubox}

\commentary
When the November contract begins to rise from its floor, a
hundred-tawara advance is a matter of course; therefore, at a mere fifty
or sixty tawara, one should not take profits. At a hundred-tawara rise,
that is the standard ceiling price, and one may take profits. However, in
a spring or summer market, if the spread against other regions is fully
developed, even a hundred-tawara rise does not necessarily mark the
ceiling --- one should keep this well in mind.

% =============================================================================
%  Article 72
% =============================================================================
\article{72}{Caution in a Range-Bound Market}

\begin{sokyubox}
{\small\itshape\color{sokyubar} \jp{第七十二章\quad 通相塲用心の事}}

\medskip
In a range-bound market (\jt{通相塲}{kayoi s\=oba}), one may have sold
four or five hundred ry\=o worth and find twenty tawara or so of profit
accrued. Thinking one has hit the mark, one piles on more --- but this is
a grave misjudgment. The swings in a range-bound market amount to only
ten or twenty tawara; one should take profits quickly and not regret
having let go of the position.
\end{sokyubox}

\commentary
When one has opened a position on the view that the market is range-bound
and, say, twenty or thirty ch\=o of profit accumulates, one is tempted by
greed to pile on, forgetting the original assessment. But this is a great
error. The typical price range in a range-bound market is only twenty or
thirty ch\=o. When ten or twenty ch\=o of profit has accrued, one should
promptly take profits (\jt{利喰}{rigui}) and close the position. The
market opens every day and one can always trade again; one should not
begrudge letting go of one's position.

% =============================================================================
%  Article 73
% =============================================================================
\article{73}{The Secret Tradition of Waiting Two Days When the Urge to Trade Surges}

\begin{sokyubox}
{\small\itshape\color{sokyubar} \jp{第七十三章\quad 賣買共心進み立候節二日相待秘傳の事}}

\medskip
When one feels compelled that this rice absolutely must rise and one must
buy today --- wait two days. When the urge to sell surges because prices
absolutely must fall --- again, wait two days. This is the innermost
secret tradition (\jt{極意の秘傳}{gokui no hiden}). In all cases, when
the market reaches the ceiling price (\jt{天井直段}{tenj\=o nedan}),
watchful timing comes first. When the ceiling price appears, concentrate
the mind solely on selling. When the floor price (\jt{底直段}{soko
nedan}) appears, concentrate the mind solely on buying. Never forget this
discipline.
\end{sokyubox}

\commentary
In Sun Tzu's chapter on strategic dispositions (\textit{Gunken-hen}), it
is said: an eye that can see the sun and moon is not thereby called
sharp, and an ear that can hear thunder is not thereby called keen ---
for any person's eyes and ears can see and hear such things. Likewise,
when the market nears the ceiling, it looks to everyone as though it will
rise; and when it nears the floor, it looks to everyone as though it will
fall. One is naturally tempted to trade in the direction of popular
sentiment. But at such moments, the causes of the rise or fall have
already been fully expressed in the price --- the market stands at the
very point of reversal. Therefore, no matter how intense the urge to
trade, one should refrain for two or three days. Invariably, the market
will move against popular sentiment.

Therefore, when one senses that the ceiling price is near, one must not
forget the resolve to sell. And when one judges that the market has
reached the floor, one must not forget the resolve to buy on the dip.
This is the foremost secret tradition of the market.

% =============================================================================
%  End of Upper Volume
% =============================================================================
\begin{center}
\rule{2in}{0.4pt}\\[1em]
{\small\itshape End of the Upper Volume}\\[0.3em]
{\small \jp{宗久翁相塲全集上卷了}}
\end{center}

\part{Lower Volume (\jp{下卷})}

\chapter*{}
\vspace{-4\baselineskip}
\markboth{Lower Volume}{}

% =============================================================================
%  Lower Volume — Article 1
% =============================================================================
\article{1}{On Margin Calls}

\begin{sokyubox}
{\small\itshape\color{sokyubar} \jp{第一章\quad 落掛りの事}}

\medskip
When one has bought and the price has risen so many tawara, the margin
call (\jt{落掛り}{ochikakari}) is triggered at such-and-such a level.

When one has sold and the price has fallen so many tawara, the margin
call is triggered at such-and-such a level.

When executing trades, one must be vigilant during periods of market
fluctuation, and manage margin calls (\jt{掛引}{kakehiki}) according to
the following rules.

When a full margin call (\jt{落切}{ochikiri}) is reached, one should
promptly buy back at the same price level. However, if the counterparty
carries the position (\jt{繋ぎ}{tsunagi}), there is no need to buy.
\end{sokyubox}

\commentary
The source includes a schematic diagram of the margin-call
(\jt{落掛り}{ochikakari}) system, showing a scale from one to ten with
corresponding levels for: \textbf{torikata} (\jp{取方}) --- the profit
side (the party whose position is in the money); \textbf{yarikata}
(\jp{遣り方}) --- the loss side (the party whose position is against
them); the margin-call trigger point (\jp{落}); and the reverse
margin-call level.

The \textit{yarikata} (loss side) instructions state: since this is a
closing-out and sell-back, even if a margin-call notice has been given,
there is no need to rush. The following day, observe the settlement and
market conditions; if the price is roughly at the same level
(\textit{ch\=o-ch\=o kurai}), one should carry the position
(\textit{tsunagi}). If one observes the market conditions carefully and the
price is roughly level, one should carry. If the position falls one
tawara further beyond the trigger, one should close out entirely and wait
to sell again when prices reach a higher level.

% =============================================================================
%  Lower Volume — Article 2 (heading only; body begins on next page)
% =============================================================================
\article{2}{When the Market Consolidates for Two or Three Months}

\begin{sokyubox}
{\small\itshape\color{sokyubar} \jp{第二章\quad 相塲二三ケ月持合ふ時心得の事}}

\medskip
When the market consolidates (\jt{保合}{mochai}) for two or three
months, eight or nine out of every ten traders turn to the sell
side. Yet afterward, the market invariably rises. The exception:
when rice has climbed relentlessly and consolidates at a level a
hundred tawara or more above the floor, one should regard it as a
falling market.
\end{sokyubox}

\commentary
As a rule, when the market enters a consolidation, sentiment turns
gloomy. Market feeling drifts toward the bearish side for no clear
reason, and a general selling mood takes hold. Yet afterward, prices
move contrary to that sentiment and rise. The reason is this:
holding-pattern markets are, in eight or nine cases out of ten,
consolidations at the floor. Everyone has sold into the
consolidation, pounding the market down, and even though it looks as
though it should fall further, it refuses to do so. If it cannot go
down, it has nowhere to go but up. The market has been storing
latent upward energy, and so, given even a small catalyst, it
erupts upward. Even so, one must assess the price level
(\jt{位取}{kurai-dori}): examining the footprint chart
(\jt{足取表}{ashitori-hy\=o}), if the market has climbed steadily
from the floor and consolidates at a level a hundred tawara or more
above it, that is a falling-market case---the exception to the
holding-pattern principle explained here.
